\chapter{Key Exchange and Forward Secrecy}

Diffie and Hellman gave us a small miracle: two strangers shouting across a
crowded room can agree on a shared secret that the room never learns. But the
plain handshake from the public-key chapter has a hole in it that you could
drive a truck through, and the truck has a name. It is called Eve, she sits in
the middle, and she is happy to be \emph{both} ends of your conversation at
once.

This chapter is about turning a raw key-agreement primitive into something you
can actually deploy. We need three things the bare handshake lacks. First,
\emph{authentication}: each party must know who is on the other end, so that the
secret is shared with Bob and not with Eve-pretending-to-be-Bob. Second,
\emph{forward secrecy}: if Bob's long-term key is stolen next year, the messages
he sent this morning must stay unreadable. Third, a disciplined way to turn the
one shared group element the handshake produces into the several independent
keys a real protocol needs --- that is the job of a \emph{key derivation
function}. Along the way we will worry about replay and downgrade, two failure
modes that have sunk real protocols. By the end you will have hand-assembled the
core of what TLS 1.3 and the Signal protocol actually do.

\section{The threat we are defending against}

Let us fix the cast and the rules before we build anything. Alice and Bob want
to establish a shared secret key over a network. Eve sits on the wire. We grant
Eve quite a lot of power, because real adversaries have it:

\begin{itemize}
  \item She sees every byte that crosses the network (eavesdropping).
  \item She can drop, reorder, replay, and modify messages, and inject her own
        (an \emph{active}, or \emph{man-in-the-middle}, adversary).
  \item She knows the protocol completely. The only things she does not know are
        the secret keys and the random coins the parties flip.
\end{itemize}

Against this adversary we want a protocol with two output guarantees.

\begin{definition}[Informal goals of authenticated key exchange]
\label{ch23:ake-goals}
A \emph{key-exchange protocol} run between two parties produces a session key
$k$. We say the protocol is a secure \emph{authenticated key exchange} (AKE) if:
\begin{enumerate}
  \item \textbf{Secrecy.} An adversary who watches (and tampers with) the
        protocol cannot distinguish the real session key $k$ from a uniformly
        random string of the same length. Concretely, $k$ should look like
        fresh randomness to anyone but the two participants.
  \item \textbf{Authentication.} When Alice finishes a run believing she shares
        $k$ with Bob, then either she really does share $k$ with the genuine Bob,
        or Bob has no matching session at all (he never accepted). She is never
        fooled into sharing a key with Eve while thinking it is Bob.
\end{enumerate}
\end{definition}

The secrecy notion above is the ``real-or-random'' game you have met for
encryption, lifted to keys. We will make it sharper when we tie security to the
Decisional Diffie-Hellman assumption.

\begin{intuitionbox}
Think of a key exchange as a blind handshake in a dark room. Secrecy says the
secret grip you settle on is meaningless to anyone groping around the edges.
Authentication says you are gripping the hand you think you are, not an
imposter's. The plain Diffie-Hellman handshake delivers the first and completely
abandons the second --- which is exactly the gap Eve lives in.
\end{intuitionbox}

\section{Recalling the bare handshake, and watching it fail}

Fix a cyclic group $\G$ of prime order $q$ with generator $g$. (Concretely think
of $\Zns{p}$ for a large prime $p$, or an elliptic-curve group; the algebra is
identical.) The textbook Diffie-Hellman exchange is one round:

\begin{center}
\begin{tikzpicture}[font=\small,>=Stealth]
  \node[draw,rounded corners,fill=aliceblue!12,minimum width=2.4cm] (A) at (0,0) {Alice};
  \node[draw,rounded corners,fill=bobgreen!12,minimum width=2.4cm] (B) at (8,0) {Bob};
  \node[below=0.1cm of A,align=center] {picks $a\getsr\Z_q$\\sends $A=g^a$};
  \node[below=0.1cm of B,align=center] {picks $b\getsr\Z_q$\\sends $B=g^b$};
  \draw[->] (1.3,0.25) -- node[above]{$A=g^a$} (6.7,0.25);
  \draw[<-] (1.3,-0.25) -- node[below]{$B=g^b$} (6.7,-0.25);
  \node[below=1.5cm of A,align=center] {computes\\$B^a=g^{ab}$};
  \node[below=1.5cm of B,align=center] {computes\\$A^b=g^{ab}$};
\end{tikzpicture}
\end{center}

Both sides arrive at $g^{ab}$. Eve sees $g^a$ and $g^b$ and would have to compute
$g^{ab}$ from them --- the \emph{Computational Diffie-Hellman} problem, believed
hard. So against a \emph{passive} Eve the shared point is secret. Good.

Now make Eve active. She does not need to solve any hard problem; she just plays
two handshakes at once.

\begin{attackbox}[Man-in-the-middle on unauthenticated DH]
\label{ch23:mitm}
Eve intercepts Alice's $A=g^a$ and never lets it reach Bob. She picks her own
$e\getsr\Z_q$ and:
\begin{itemize}
  \item sends $E=g^e$ to Bob, pretending to be Alice;
  \item sends $E=g^e$ to Alice, pretending to be Bob (she can reuse $e$ or pick a
        fresh one --- say $e'$ toward Alice).
\end{itemize}
Now Alice computes $k_A=E^a=g^{ae}$ and thinks she shares it with Bob. Bob
computes $k_B=E^b=g^{be}$ and thinks he shares it with Alice. Eve knows $e$, so
she computes \emph{both}: $k_A=A^e$ and $k_B=B^e$. She decrypts everything Alice
sends under $k_A$, re-encrypts it under $k_B$, and forwards it to Bob. Neither
party sees anything wrong. Eve is a transparent relay who reads and rewrites the
whole conversation.
\end{attackbox}

The lesson is blunt: \textbf{a key is only as trustworthy as your knowledge of
who holds the other copy.} Diffie-Hellman gives you a shared secret with
\emph{somebody}. Turning ``somebody'' into ``Bob'' is the entire problem of
authentication, and it cannot come from the group algebra alone. It must come
from something Alice already knows about Bob: a long-term key, bound to his
identity, that Eve cannot forge.

\section{Authenticating the handshake}

We have two long-term-key tools that can pin an identity: digital signatures and
long-term Diffie-Hellman keys. They lead to two families of AKE, and modern
protocols use both. Let us build the signature-based one first because it is the
most transparent, then the implicitly-authenticated one because it is what
elliptic-curve protocols and Signal actually run.

\subsection{Signed ephemeral Diffie-Hellman}

Give Bob a long-term signature key pair $(sk_B, pk_B)$, and suppose Alice already
holds an authentic copy of $pk_B$ --- through a certificate, a pinned key, or a
trust-on-first-use record. (Where $pk_B$ comes from is the job of the PKI
chapter; here we just assume Alice has the right one.) The idea is to keep the
Diffie-Hellman values \emph{ephemeral} --- fresh per session --- but have each
party \emph{sign} the ephemeral value they send.

\begin{center}
\begin{tikzpicture}[font=\small,>=Stealth]
  \node[draw,rounded corners,fill=aliceblue!12,minimum width=2.6cm] (A) at (0,0) {Alice};
  \node[draw,rounded corners,fill=bobgreen!12,minimum width=2.6cm] (B) at (9,0) {Bob};
  \draw[->] (1.4,0.55) -- node[above,align=center]{$A=g^a$ \quad (``ClientHello'')} (7.6,0.55);
  \draw[<-] (1.4,-0.1) -- node[below,align=center]{$B=g^b,\ \sigma_B=\Sign_{sk_B}(A\concat B)$} (7.6,-0.1);
  \draw[->] (1.4,-1.0) -- node[below,align=center]{(optional) $\sigma_A=\Sign_{sk_A}(A\concat B)$} (7.6,-1.0);
\end{tikzpicture}
\end{center}

\begin{enumerate}
  \item Alice picks $a\getsr\Z_q$ and sends $A=g^a$.
  \item Bob picks $b\getsr\Z_q$, and signs the \emph{transcript} so far:
        $\sigma_B=\Sign_{sk_B}(A\concat B)$. He sends $(B,\sigma_B)$.
  \item Alice checks $\Vrfy_{pk_B}(A\concat B,\sigma_B)$. If it verifies, she
        knows the genuine Bob saw \emph{her} $A$ and produced \emph{this} $B$.
        She computes $k=B^a=g^{ab}$.
  \item Bob computes $k=A^b=g^{ab}$. (If Alice also has a long-term key and Bob
        needs to authenticate her --- ``mutual authentication'' --- she signs
        the transcript too in step 3.)
\end{enumerate}

Why does signing the transcript, rather than just signing $B$, matter so much?
Because the signature must \emph{bind Bob's ephemeral to Alice's ephemeral}. If
Bob signed only $B$, Eve could record an old $(B,\sigma_B)$ pair from a previous
session and replay it into a new one --- the signature would still verify. By
covering $A\concat B$, the signature is welded to this specific run. Alice put a
fresh, random $A$ on the wire; a valid signature over $A\concat B$ proves Bob
acted \emph{after} seeing that fresh $A$, so the message cannot be a recording.

\begin{whybox}[Why MITM dies here]
Return to Eve from box~\ref{ch23:mitm}. To impersonate Bob to Alice, Eve must
send some $B'$ and a signature $\sigma'$ with $\Vrfy_{pk_B}(A\concat B',\sigma')$
accepting. But producing a valid signature under $pk_B$ on a message of her
choosing is exactly an existential forgery --- the thing a secure signature
scheme makes infeasible. Eve does not have $sk_B$. She is stuck. She can still
relay the genuine messages unchanged, but then she has merely forwarded a real
Alice--Bob exchange and learned nothing: the resulting $g^{ab}$ is secret from
her by Computational Diffie-Hellman. The signature converts ``a key with
somebody'' into ``a key with whoever owns $sk_B$,'' and the certificate says that
owner is Bob.
\end{whybox}

\begin{trapbox}[Sign the transcript, not just your own share]
A recurring real-world bug is signing too little. If a party signs only its own
public value, or signs a message that omits the peer's nonce, an attacker can
splice signed pieces from different sessions together. The discipline ``everyone
signs a hash of the full transcript they have seen so far'' is what TLS 1.3's
\texttt{CertificateVerify} does, and it is the safe default. We will see the
matching idea --- a final MAC over the whole transcript --- in the downgrade
discussion.
\end{trapbox}

\subsection{The crucial word: ephemeral}

Notice that $a$ and $b$ are freshly chosen for this session and \emph{thrown away
afterward}. The long-term keys ($sk_B$, and $sk_A$ if used) sign, but they never
touch the shared secret $g^{ab}$. That separation --- long-term keys for
\emph{authentication}, ephemeral keys for \emph{secrecy} --- is the structural
move that buys forward secrecy. Hold that thought; it is the heart of the next
section.

\begin{historybox}
The ``station-to-station'' (STS) protocol of Diffie, van Oorschot, and Wiener
(1992) is essentially the signed-ephemeral-DH scheme above, with a refinement:
each party signs the DH values and then MACs that signature under the new
session key, which proves they actually computed the same key (``key
confirmation''). STS is the direct ancestor of the authenticated handshakes in
IKE (the IPsec key-exchange protocol), SSH, and TLS. The lineage from a 1976
idea to the handshake protecting your bank login runs straight through here.
\end{historybox}

\subsection{Implicit authentication: long-term DH keys}

Signatures are not the only way. Suppose instead Bob's long-term key \emph{is} a
Diffie-Hellman key: a secret $b_{\mathrm{lt}}$ with public value
$B_{\mathrm{lt}}=g^{b_{\mathrm{lt}}}$ that Alice trusts. Alice similarly has
$a_{\mathrm{lt}}$, $A_{\mathrm{lt}}=g^{a_{\mathrm{lt}}}$. Then Alice and Bob can
agree on a key with no signatures at all, because only the true owners of the
long-term secrets can compute the right shared value.

The cleanest illustration is the \emph{triple Diffie-Hellman} pattern used by
Signal. Each party has a long-term ``identity'' key and a fresh ephemeral key:

\[
\text{Alice: } (a_{\mathrm{lt}},\, a_{\mathrm{eph}}),\qquad
\text{Bob: } (b_{\mathrm{lt}},\, b_{\mathrm{eph}}),
\]
with the corresponding public values exchanged. They form the shared secret by
hashing together several DH outputs:
\[
k = \Hash\!\big(\;
g^{a_{\mathrm{lt}}\, b_{\mathrm{eph}}}\concat
g^{a_{\mathrm{eph}}\, b_{\mathrm{lt}}}\concat
g^{a_{\mathrm{eph}}\, b_{\mathrm{eph}}}
\;\big).
\]
\begin{figure}[H]
\centering
\begin{tikzpicture}[font=\small,>=Stealth,node distance=6mm]
  \node[align=center,fill=aliceblue!12,rounded corners,draw,minimum width=2.0cm] (alt) at (0,1.4) {$a_{\mathrm{lt}}$\\(Alice long-term)};
  \node[align=center,fill=aliceblue!12,rounded corners,draw,minimum width=2.0cm] (aeph) at (0,-1.4) {$a_{\mathrm{eph}}$\\(Alice ephemeral)};
  \node[align=center,fill=bobgreen!12,rounded corners,draw,minimum width=2.0cm] (blt) at (6,1.4) {$b_{\mathrm{lt}}$\\(Bob long-term)};
  \node[align=center,fill=bobgreen!12,rounded corners,draw,minimum width=2.0cm] (beph) at (6,-1.4) {$b_{\mathrm{eph}}$\\(Bob ephemeral)};
  \draw[->,thick] (alt) -- node[pos=0.26,above,sloped,font=\scriptsize]{authenticates Alice} (beph);
  \draw[->,thick] (aeph) -- node[pos=0.26,below,sloped,font=\scriptsize]{authenticates Bob} (blt);
  \draw[->,thick,structurepurple] (aeph) -- node[below,font=\scriptsize]{fresh randomness} (beph);
\end{tikzpicture}
\caption{The three Diffie-Hellman cross-products of triple DH. Each long-term-against-ephemeral
arrow pins one party's identity; the ephemeral-against-ephemeral arrow injects per-session
freshness for forward secrecy.}
\label{ch23:tripledh}
\end{figure}

Read each factor as a claim. The first mixes Alice's \emph{long-term} key with
Bob's \emph{ephemeral} --- only the real Alice (who knows $a_{\mathrm{lt}}$) and
the sender of $b_{\mathrm{eph}}$ can form it, so it authenticates Alice. The
second mixes Bob's long-term with Alice's ephemeral --- it authenticates Bob.
The third mixes the two ephemerals --- it injects fresh, throwaway randomness so
that the session secret depends on per-session coins (Figure~\ref{ch23:tripledh}). ``Implicit'' means there
is no separate signature to verify: a party who can compute the same $k$ has, by
that fact alone, proven it holds the right long-term secret.

\begin{intuitionbox}
A signature is a notarized letter: ``I, Bob, hereby endorse this handshake,''
verifiable by anyone with Bob's public key. Implicit authentication is more like
a secret handshake whose grip only works if both palms carry the right tattoo.
Nobody signs anything; the key simply will not come out right unless you are who
you claim to be. Signal prefers the implicit style partly because it leaves no
signature on the wire --- a small but real privacy gain, since a signature is a
non-repudiable record that you said a particular thing.
\end{intuitionbox}

\section{Forward secrecy}

Here is the scenario that forward secrecy is built to survive. Eve cannot break
your handshake today, so she does the patient thing: she records every encrypted
session you have, byte for byte, and stores them. Months later she compromises
your server --- a stolen laptop, a subpoena, a Heartbleed-style memory leak ---
and walks away with your \emph{long-term private key}. The question that decides
whether this is a catastrophe or a shrug:

\begin{center}
\emph{Does possessing the long-term key today let Eve decrypt the traffic she
recorded yesterday?}
\end{center}

\begin{definition}[Perfect forward secrecy]
\label{ch23:pfs}
A key-exchange protocol has \emph{(perfect) forward secrecy} if the compromise of
long-term secret keys does not compromise the session keys of \emph{past}
sessions. Equivalently: an adversary who records a session and \emph{later}
learns all the long-term private keys still cannot recover that session's key.
\end{definition}

The word ``perfect'' here is historical and a bit grandiose; people increasingly
just say ``forward secrecy.'' The property is entirely about a \emph{time
ordering}: secrets created and destroyed before the compromise stay safe.

\subsection{Why ephemeral DH gives it, and static encryption does not}

Picture two ways Alice could send Bob a session key.

\textbf{Way 1 --- static key transport (no forward secrecy).} Alice picks a
random session key $k$, encrypts it under Bob's long-term public key, and sends
$\Enc_{pk_B}(k)$. This is essentially old TLS with RSA key transport. Eve records
the ciphertext. Years later she steals $sk_B$. She computes
$\Dec_{sk_B}(\Enc_{pk_B}(k)) = k$ and decrypts the whole recorded session. The
long-term key is a \emph{master key to the past}. One compromise unlocks every
session ever recorded under it. This is a disaster, and it is why this mode was
removed from TLS 1.3.

\textbf{Way 2 --- ephemeral DH (forward secrecy).} Alice and Bob run the signed
ephemeral handshake. The session key is $g^{ab}$, derived from the ephemeral
secrets $a$ and $b$, which both parties \emph{erase} as soon as the handshake
finishes. The long-term key only \emph{signed}; it never touched $g^{ab}$. Now
Eve steals $sk_B$ later. What can she do with it? She can forge signatures from
\emph{now on} --- she can impersonate Bob in \emph{future} handshakes. But for
the recorded past session she has only $g^a$, $g^b$, and a signature. To recover
$g^{ab}$ from $g^a$ and $g^b$ is, again, Computational Diffie-Hellman --- and the
stolen signing key is no help at all, because it has nothing to do with the DH
secret. The ephemerals $a,b$ are gone, deleted. The past stays dark.

\begin{whybox}[The one-sentence reason]
Forward secrecy works because the session secret is computed from
\emph{ephemeral} values that are destroyed after use, while the long-term key
only \emph{authenticates} and never enters the secret. Compromising the
authenticator cannot retroactively reconstruct secrets it never saw.
\end{whybox}

\begin{figure}[H]
\centering
\begin{tikzpicture}[font=\small, >=stealth]
  % Lane 1: Static Key Transport (No Forward Secrecy)
  \begin{scope}[shift={(0,4.2)}]
    \node[anchor=south west, font=\footnotesize\scshape, text=trapred!90!black] at (0,3.3) {Scenario A: Static Key Transport (No Forward Secrecy)};
    \draw[draw=trapred!40, fill=trapred!5, rounded corners=2pt] (0,0) rectangle (14.5,3.2);

    % Time axis
    \draw[->, thick, gray!70] (0.8,0.5) -- (14.0,0.5) node[right, font=\scriptsize, text=black] {Time $t$};
    
    % Session 1
    \node[draw=sourceblue!60, fill=sourceblue!15, rounded corners=2pt, font=\scriptsize, align=center] (s1) at (2.4,1.4) {
      \textbf{Session 1}\\$c_1 = \Enc_{pk_B}(k_1)$
    };
    \draw[thick, gray!50] (2.4,0.5) -- (s1.south);
    \node[font=\tiny, text=gray!80!black, below] at (2.4,0.45) {$t_1$};

    % Session 2
    \node[draw=sourceblue!60, fill=sourceblue!15, rounded corners=2pt, font=\scriptsize, align=center] (s2) at (5.2,1.4) {
      \textbf{Session 2}\\$c_2 = \Enc_{pk_B}(k_2)$
    };
    \draw[thick, gray!50] (5.2,0.5) -- (s2.south);
    \node[font=\tiny, text=gray!80!black, below] at (5.2,0.45) {$t_2$};

    % Session 3
    \node[draw=sourceblue!60, fill=sourceblue!15, rounded corners=2pt, font=\scriptsize, align=center] (s3) at (8.0,1.4) {
      \textbf{Session 3}\\$c_3 = \Enc_{pk_B}(k_3)$
    };
    \draw[thick, gray!50] (8.0,0.5) -- (s3.south);
    \node[font=\tiny, text=gray!80!black, below] at (8.0,0.45) {$t_3$};

    % Compromise Point
    \node[draw=trapred, fill=trapred!20, rounded corners=2pt, font=\scriptsize\bfseries, align=center, text=trapred!90!black] (comp1) at (11.8,1.4) {
      Key Compromise!\\Steal $sk_B$
    };
    \draw[thick, trapred] (11.8,0.5) -- (comp1.south);
    \node[font=\tiny, text=trapred!90!black, below] at (11.8,0.45) {$T_{\mathrm{compromise}}$};

    % Retroactive decryption arc above the boxes
    \draw[->, very thick, trapred, dashed] (comp1.north) to[out=140, in=40] node[above=2pt, font=\scriptsize\bfseries, text=trapred]{Retroactive Decryption: $k_i = \Dec_{sk_B}(c_i)$ (all past recorded traffic exposed!)} (s1.north);
  \end{scope}

  % Lane 2: Ephemeral Diffie-Hellman (Perfect Forward Secrecy)
  \begin{scope}[shift={(0,0)}]
    \node[anchor=south west, font=\footnotesize\scshape, text=targetgreen!90!black] at (0,3.3) {Scenario B: Ephemeral Diffie--Hellman (ECDHE with Perfect Forward Secrecy)};
    \draw[draw=targetgreen!40, fill=targetgreen!5, rounded corners=2pt] (0,0) rectangle (14.5,3.2);

    % Time axis
    \draw[->, thick, gray!70] (0.8,0.5) -- (14.0,0.5) node[right, font=\scriptsize, text=black] {Time $t$};

    % Session 1
    \node[draw=targetgreen!60, fill=targetgreen!15, rounded corners=2pt, font=\scriptsize, align=center] (es1) at (2.4,1.4) {
      \textbf{Session 1: } $g^{a_1 b_1}$\\
      \tiny $a_1, b_1$ {\color{evered}\textbf{erased}}
    };
    \draw[thick, gray!50] (2.4,0.5) -- (es1.south);
    \node[font=\tiny, text=gray!80!black, below] at (2.4,0.45) {$t_1$};

    % Session 2
    \node[draw=targetgreen!60, fill=targetgreen!15, rounded corners=2pt, font=\scriptsize, align=center] (es2) at (5.2,1.4) {
      \textbf{Session 2: } $g^{a_2 b_2}$\\
      \tiny $a_2, b_2$ {\color{evered}\textbf{erased}}
    };
    \draw[thick, gray!50] (5.2,0.5) -- (es2.south);
    \node[font=\tiny, text=gray!80!black, below] at (5.2,0.45) {$t_2$};

    % Session 3
    \node[draw=targetgreen!60, fill=targetgreen!15, rounded corners=2pt, font=\scriptsize, align=center] (es3) at (8.0,1.4) {
      \textbf{Session 3: } $g^{a_3 b_3}$\\
      \tiny $a_3, b_3$ {\color{evered}\textbf{erased}}
    };
    \draw[thick, gray!50] (8.0,0.5) -- (es3.south);
    \node[font=\tiny, text=gray!80!black, below] at (8.0,0.45) {$t_3$};

    % Compromise Point
    \node[draw=targetgreen!80!black, fill=targetgreen!20, rounded corners=2pt, font=\scriptsize\bfseries, align=center, text=black] (comp2) at (11.8,1.4) {
      Key Compromise!\\Steal $sk_B^{\mathrm{sign}}$
    };
    \draw[thick, trapred] (11.8,0.5) -- (comp2.south);
    \node[font=\tiny, text=trapred!90!black, below] at (11.8,0.45) {$T_{\mathrm{compromise}}$};

    % Blocked retroactive decryption banner in clear corridor
    \node[draw=targetgreen!80!black, fill=targetgreen!20, rounded corners=2pt, font=\scriptsize\bfseries, text=targetgreen!90!black] at (5.6,2.5) {
      $\times$ Past Sessions Protected: ephemeral keys were zeroized; recorded wires stay dark!
    };
  \end{scope}
\end{tikzpicture}
\caption{Timeline comparison of session key security upon long-term key compromise.
In static key transport (Scenario A), the long-term secret key acts as a master key to all past recorded traffic.
In ephemeral Diffie--Hellman (Scenario B), past session keys are derived from transient exponents zeroized
after the handshake; compromise of the long-term signing key allows only future impersonation,
never retroactive decryption of the past.}
\label{ch23:fig-pfs-timeline}
\end{figure}

\begin{trapbox}[``Ephemeral'' is only as good as your delete]
Forward secrecy is a property of the \emph{whole system}, not just the math. If
your server keeps the ephemeral secret $b$ in memory, or worse logs it, or reuses
one ``ephemeral'' value across thousands of connections to save CPU, then
stealing it resurrects all those sessions. There is a real-world cautionary tale
here: some servers cached a single DH ephemeral for the lifetime of the process
(sometimes called a static-ephemeral or ``semi-static'' key) to avoid recomputing
$g^b$. That quietly destroys forward secrecy for every session sharing the cached
value. Ephemeral means \emph{generated fresh and wiped promptly}. If you cannot
guarantee the wipe, you do not have the property.
\end{trapbox}

\begin{applicationbox}
This is why every modern TLS deployment insists on ``(EC)DHE'' cipher suites ---
the trailing \texttt{E} is for \emph{ephemeral}. TLS 1.3 went further and made
ephemeral (elliptic-curve) Diffie-Hellman the \emph{only} key-establishment
mechanism: static RSA key transport was deleted from the protocol outright,
precisely because it lacked forward secrecy and made mass retroactive decryption
a one-key-theft-away proposition. If you operate a server, ``do we have forward
secrecy'' reduces to ``are we negotiating an ephemeral DH suite and actually
discarding the ephemerals.''
\end{applicationbox}

\subsection{A worked tiny example}

Numbers make the time-ordering concrete. Work in $\Zns{p}$ with the small prime
$p=23$. The element $g=5$ generates the group ($5$ has order $22$ modulo $23$, so
it is a generator). All arithmetic is mod $23$.

\textbf{The session.} Alice picks ephemeral $a=6$, sends $A = 5^6 \bmod 23$.
Compute: $5^2=25\equiv 2$, $5^3\equiv 10$, $5^6=(5^3)^2\equiv 10^2=100\equiv 8$.
So $A=8$.

Bob picks ephemeral $b=15$, sends $B = 5^{15}\bmod 23$. Compute by squaring:
$5^1=5$, $5^2=2$, $5^4=2^2=4$, $5^8=4^2=16$. Then
$5^{15}=5^{8}\cdot5^{4}\cdot5^{2}\cdot5^{1}=16\cdot4\cdot2\cdot5$. Step it out:
$16\cdot4=64\equiv 18$; $18\cdot2=36\equiv 13$; $13\cdot5=65\equiv 19$. So $B=19$.

\textbf{Shared secret.} Alice computes $B^a=19^{6}\bmod 23$ and Bob computes
$A^b=8^{15}\bmod 23$; both equal $g^{ab}=5^{90}\bmod 23$. Since the group has
order $22$, reduce the exponent: $90 = 4\cdot 22 + 2$, so
$5^{90}\equiv 5^{2}\equiv 2 \pmod{23}$. The session secret is $2$. (You can check
$19^6$ directly: $19\equiv -4$, $(-4)^6=4096$, and $4096 = 178\cdot 23 + 2$, so
$19^6\equiv 2$. Agreement.)

\textbf{Now the theft.} Suppose tomorrow Eve steals Bob's long-term \emph{signing}
key. From the recorded wire she has $A=8$ and $B=19$. The signing key lets her
forge signatures but tells her nothing about $a=6$ or $b=15$, which Alice and Bob
deleted. To get the secret $2$ she must solve $5^x\equiv 8$ (find $a$) or
$5^y\equiv 19$ (find $b$) --- a discrete log. With $p=23$ she could, of course,
by brute force; that is the price of toy numbers. With a real $2048$-bit or
$256$-bit-curve group, that search is the wall the whole field is built on. The
structural point survives the small numbers: the stolen long-term key is no
shortcut to the recorded session's secret.

\section{From one group element to many keys: KDFs}

The handshake hands you a single shared secret --- a group element like our $2$,
or in practice a $256$-bit string. But a real session needs \emph{several}
independent keys: one to encrypt client-to-server, one for server-to-client,
maybe separate MAC keys, IVs, a key to confirm the handshake. You must not just
slice the raw DH output into chunks and use them directly. Two reasons.

First, the raw shared secret is not a uniform random string. A Diffie-Hellman
value $g^{ab}$ in $\Zns{p}$ is an element of a structured set --- it is biased,
it has algebraic relations, its bits are not independent. Cryptographic keys are
supposed to be (indistinguishable from) uniform. We need to ``iron out'' the
structure.

Second, you want \emph{domain separation}: keys for different purposes must be
provably unrelated, so that a flaw in one use cannot leak another. Reusing one
key for two jobs is a classic footgun.

The tool for both is a \emph{key derivation function} (KDF). The standard,
HKDF, splits the job into two stages with the slogan \emph{extract-then-expand}.

\begin{definition}[HKDF, extract-then-expand]
\label{ch23:hkdf}
Built from an HMAC keyed by a hash $\Hash$:
\begin{itemize}
  \item \textbf{Extract.} From the raw, non-uniform input keying material
        $\mathrm{IKM}$ (here $\mathrm{IKM}=g^{ab}$) and an optional non-secret
        \emph{salt}, produce a fixed-length, uniform-looking pseudorandom key:
        \[ \mathrm{PRK} \;=\; \mathrm{HMAC}(\text{salt},\ \mathrm{IKM}). \]
        This is the ``ironing'' step. HMAC, used as a randomness extractor,
        concentrates whatever entropy is in $\mathrm{IKM}$ into a full-strength
        uniform key $\mathrm{PRK}$.
  \item \textbf{Expand.} From $\mathrm{PRK}$ and a per-key \emph{info} label,
        stretch out as many key bytes as you need, one HMAC block at a time:
        \[ T_i = \mathrm{HMAC}\big(\mathrm{PRK},\ T_{i-1}\concat \mathrm{info}\concat i\big),
           \quad T_0 = \text{empty}, \]
        and concatenate $T_1\concat T_2\concat\cdots$ to the requested length.
        The \emph{info} string is where domain separation lives: feed
        \texttt{"c hs traffic"} for one key and \texttt{"s hs traffic"} for
        another, and the two outputs are unrelated.
\end{itemize}
\end{definition}

\begin{intuitionbox}
\emph{Extract} is a juicer: the raw DH element is lumpy fruit with entropy
unevenly distributed; HMAC squeezes out a clean, uniform concentrate
($\mathrm{PRK}$). \emph{Expand} is a vending machine: insert $\mathrm{PRK}$,
punch a label (\texttt{"enc"}, \texttt{"mac"}, \texttt{"iv"}), and out comes a
fresh independent key. Different labels, different keys, with a security
guarantee that knowing one tells you nothing about the others. One handshake
secret, a whole keyring out the other side.
\end{intuitionbox}

\begin{connectionbox}
Why two stages instead of just ``hash the DH output''? It is the same
distinction the probability companion draws between an \emph{extractor} (turn a
weak, biased source into near-uniform bits --- a min-entropy and leftover-hash
argument) and a \emph{pseudorandom generator} (stretch a short uniform seed into
a long pseudorandom stream). Extract handles the non-uniform \emph{source};
expand handles the \emph{length}. Conflating them is how subtle biases sneak from
the group into your keys. The HMAC indistinguishability that makes expand safe is
the same PRF property developed in the MAC chapter.
\end{connectionbox}

\begin{implementbox}[Toy HKDF over a real hash --- learning only]
This shows the extract-then-expand shape. It is a teaching toy: no constant-time
care, no memory wiping, and it leans on the standard library for HMAC. \textbf{Do
not use this in production --- use your platform's vetted HKDF.}
\begin{lstlisting}[language=Python]
import hmac, hashlib

def _hmac(key, msg):
    return hmac.new(key, msg, hashlib.sha256).digest()

def hkdf_extract(salt, ikm):
    # salt acts as the HMAC key; if absent, use a zero block
    if not salt:
        salt = b"\x00" * hashlib.sha256().digest_size
    return _hmac(salt, ikm)            # PRK: one uniform pseudorandom key

def hkdf_expand(prk, info, length):
    out, t, counter = b"", b"", 1
    while len(out) < length:
        # T_i = HMAC(PRK, T_{i-1} || info || i)
        t = _hmac(prk, t + info + bytes([counter]))
        out += t
        counter += 1
    return out[:length]

def hkdf(salt, ikm, info, length):
    return hkdf_expand(hkdf_extract(salt, ikm), info, length)

# One shared DH secret -> several independent keys via different labels
shared = b"raw-diffie-hellman-output-not-uniform"
prk_salt = b"handshake transcript hash goes here"
enc_key = hkdf(prk_salt, shared, b"c2s enc", 32)   # client->server
dec_key = hkdf(prk_salt, shared, b"s2c enc", 32)   # server->client
mac_key = hkdf(prk_salt, shared, b"c2s mac", 32)
print(enc_key.hex())
print(dec_key.hex())     # unrelated to enc_key, by domain separation
\end{lstlisting}
The three keys come from one shared secret yet are computationally independent:
flipping the \texttt{info} label is all it takes. In TLS 1.3 the ``salt'' is
literally a running hash of the handshake transcript, which is how the derived
keys also become bound to everything said so far --- a point we exploit next.
\end{implementbox}

\section{Replay and downgrade}

Two active attacks deserve their own treatment because they bypass secrecy
entirely and attack the \emph{protocol logic}.

\subsection{Replay}

A \emph{replay} attack records a valid protocol message and re-sends it later.
The classic damage: Eve records Alice's encrypted ``transfer \$100 to Eve''
message and replays it ten times. Each copy is a perfectly valid ciphertext;
encryption alone does not notice it is a duplicate.

The defense in key exchange is \emph{freshness}: weave an unpredictable,
single-use value into each session so that a recording from one run is invalid in
another. Ephemeral DH already does this for you --- the secret is $g^{ab}$ with
fresh $a,b$, so a captured handshake cannot be replayed to produce a usable key
in a new session, because the peer chose a new ephemeral. More generally,
protocols add explicit \emph{nonces} (numbers used once) or per-session
\emph{counters}, and require every authenticated message to cover them.

\begin{trapbox}[The 0-RTT replay hole in TLS 1.3]
TLS 1.3 offers a ``0-RTT'' (zero round-trip) mode where a returning client can
send application data in its very first flight, encrypted under a key derived
from a secret remembered from a previous connection. The speed is wonderful and
the security caveat is sharp: because that first-flight data is \emph{not} bound
to a fresh server ephemeral (the server has not spoken yet), it is
\emph{replayable}. Eve can capture the 0-RTT data and re-send it to the server
multiple times. The protocol therefore restricts 0-RTT to \emph{idempotent}
requests (a \texttt{GET} you can safely repeat), never a ``transfer money''
\texttt{POST}, and servers maintain anti-replay caches. Freshness lost is replay
gained --- the trade is explicit and it is the application's job to respect it.
\end{trapbox}

\subsection{Downgrade}

A \emph{downgrade} attack does not break strong crypto; it tricks the parties
into \emph{negotiating weak crypto}. Real protocols support many cipher suites
and versions for backward compatibility. Negotiation usually starts with the
client advertising what it supports in cleartext --- before any key exists to
protect that list. So Eve, sitting in the middle, edits the list.

\begin{attackbox}[Strip the strong options]
Alice's ClientHello offers, say, modern ephemeral-ECDH suites \emph{and}, for an
ancient peer, an export-grade $512$-bit DH suite. Eve intercepts it and deletes
every strong option, leaving only the weak one. Bob, seeing only the weak choice,
agrees to it. Now the session uses $512$-bit DH, which Eve can break with modest
precomputation --- and she reads everything. Neither party intended weak crypto;
the negotiation itself was unauthenticated. This is the shape of the real
\textsc{Logjam} and \textsc{FREAK} attacks (2015), which forced browsers and
servers to drop export-grade suites entirely.
\end{attackbox}

The defense is \emph{transcript authentication}: at the end of the handshake,
once a key exists, each party computes a MAC (or signature) over a hash of
\emph{every handshake message it sent and received}, and the peer checks it.

\begin{whybox}[Why a final transcript MAC stops downgrade]
If Eve edited the negotiation, Alice's recorded transcript and Bob's recorded
transcript \emph{disagree} --- Alice ``sent'' a strong-suites list, Bob
``received'' a stripped one. When each side MACs its own transcript hash under
the freshly derived key and sends it for the other to verify, the two MACs cover
different inputs and \emph{will not match}. The handshake aborts. Eve cannot fix
this by forging a matching MAC, because the MAC key is the session secret she
does not have. So any mid-handshake tampering, including a downgrade edit, is
detected the moment the parties cross-check transcripts. This ``Finished''
message is exactly how TLS seals its handshake, and it is the same key-confirmation
idea that the station-to-station protocol introduced.
\end{whybox}

\begin{intuitionbox}
Downgrade is a forged restaurant menu: Eve hands Bob a menu with only the
spoiled dish listed, hoping he orders it. The transcript MAC is Alice and Bob, at
the end of the meal, reading their menus aloud to each other. If the menus
differ, somebody tampered, and they walk out. The negotiation can happen in the
clear; it just has to be \emph{ratified} under the key once the key exists.
\end{intuitionbox}

\section{Putting it together: the shape of TLS 1.3 and Signal}

We now have every piece. Let us see the two protocols that assemble them, because
they make different choices for different reasons.

\subsection{TLS 1.3 in one breath}

\begin{center}
\begin{tikzpicture}[font=\footnotesize,>=Stealth]
  \node[draw,rounded corners,fill=aliceblue!12] (C) at (0,0) {Client};
  \node[draw,rounded corners,fill=bobgreen!12] (S) at (10,0) {Server};
  \draw[->] (1.0,0.9) -- node[above,align=center]{ClientHello: nonce, supported suites,\\ ephemeral $g^a$} (9.0,0.9);
  \draw[<-] (1.0,0.25) -- node[below,align=center]{ServerHello: nonce, chosen suite, ephemeral $g^b$} (9.0,0.25);
  \draw[<-] (1.0,-0.55) -- node[below,align=center]{\{Certificate, CertificateVerify=$\Sign_{sk_S}(\text{transcript})$,\\ Finished=MAC(transcript)\}} (9.0,-0.55);
  \draw[->] (1.0,-1.65) -- node[below,align=center]{\{Finished=MAC(transcript)\}, application data} (9.0,-1.65);
  \node[font=\scriptsize,gray] at (5,-2.3) {\{$\cdots$\} = encrypted under handshake keys derived by HKDF from $g^{ab}$};
\end{tikzpicture}
\end{center}

Trace the parts you built:
\begin{itemize}
  \item The exchanged $g^a,g^b$ are \emph{ephemeral} ECDH --- forward secrecy by
        construction (Definition~\ref{ch23:pfs}).
  \item \texttt{CertificateVerify} is the server \emph{signing the transcript}
        with its long-term key --- signed ephemeral DH, authenticating the server
        (and the certificate ties $pk_S$ to the hostname, the PKI chapter's job).
  \item The handshake and application keys come from $g^{ab}$ through HKDF, with
        the transcript hash as salt --- one secret, a keyring, all bound to what
        was said.
  \item The two \texttt{Finished} MACs are the \emph{transcript authentication}
        that kills downgrade and confirms both sides derived the same key.
  \item Nonces in the Hellos provide \emph{freshness} against replay.
\end{itemize}
Every defense in this chapter appears exactly once. That is not a coincidence;
TLS 1.3 was designed by stripping the protocol down to precisely these
guarantees and removing everything (static RSA, renegotiation, weak suites) that
violated them.

\subsection{Signal in one breath}

Signal solves a harder problem: the recipient is usually \emph{offline} when the
first message arrives, so there is no live round trip. Its X3DH (``extended
triple Diffie-Hellman'') handshake lets Alice start a session against a bundle of
Bob's public keys that he published to a server in advance.

Bob pre-publishes an \emph{identity key} (long-term), a \emph{signed prekey}
(medium-term, signed by his identity key so Alice can trust it), and a batch of
one-time \emph{prekeys}. Alice fetches a bundle, generates her own ephemeral, and
computes the multi-DH secret from the earlier section --- mixing her ephemeral
and identity keys against Bob's signed prekey, identity key, and a one-time
prekey:
\[
k=\Hash\big(\,
\mathrm{DH}(a_{\mathrm{id}},\,B_{\mathrm{spk}})\concat
\mathrm{DH}(a_{\mathrm{eph}},\,B_{\mathrm{id}})\concat
\mathrm{DH}(a_{\mathrm{eph}},\,B_{\mathrm{spk}})\concat
\mathrm{DH}(a_{\mathrm{eph}},\,B_{\mathrm{otp}})\,\big).
\]
Read the pieces with the eyes we trained: identity-against-something terms supply
\emph{implicit authentication} of each party; the ephemeral terms supply
\emph{forward secrecy}; the one-time prekey $B_{\mathrm{otp}}$, consumed and
deleted after a single use, adds replay resistance and an extra layer of forward
secrecy for that very first message. The signed prekey carries the
\emph{signature-based} authentication (Bob's identity key signs it). Then HKDF
turns $k$ into the initial chain keys, and the ``double ratchet'' takes over from
there to refresh keys per message --- pushing forward secrecy down to the
granularity of individual texts.

\begin{applicationbox}
The contrast is instructive. TLS authenticates with a \emph{signature} over the
transcript because a web server is online and wants any client to verify it
against a public certificate chain. Signal authenticates \emph{implicitly} with
long-term DH keys because it wants asynchronous, no-round-trip session setup and
prefers to leave no signed, non-repudiable record of who said what. Both reach
the same destination --- an authenticated, forward-secret session key fed through
HKDF --- by the routes that fit their setting. WireGuard's handshake (the modern
VPN) is a third point in the same design space: it uses the implicit, DH-only
``Noise'' framework, no certificates at all, with both parties' static keys
pre-known.
\end{applicationbox}

\begin{focusbox}[The whole chapter on one page]
A deployable key exchange is ephemeral Diffie-Hellman (for \emph{forward
secrecy}) plus an \emph{authenticator} bound to the transcript (a signature, as
in TLS, or implicit long-term-DH terms, as in Signal) --- which kills the
\emph{man-in-the-middle} --- plus \emph{nonces} (against replay) and a final
\emph{transcript MAC} (against downgrade) --- with the one shared secret fanned
out into many independent keys by \emph{HKDF}. Lose the ephemerals' freshness and
you lose forward secrecy; lose the transcript binding and you lose authentication
or downgrade protection. Every real handshake is some arrangement of exactly
these moving parts.
\end{focusbox}

\section{Exercises}

\begin{exercise}
In the man-in-the-middle attack of box~\ref{ch23:mitm}, Eve must actively relay
messages for the whole session. Suppose she goes offline halfway through, after
Alice and Bob have each computed their keys $k_A=g^{ae}$ and $k_B=g^{be}$.
Describe exactly what breaks for Alice and Bob, and explain why this means a MITM
on unauthenticated DH must stay present for the entire conversation.
\end{exercise}

\begin{exercise}
Work the toy handshake in $\Zns{23}$ with $g=5$ for new ephemerals $a=9$ and
$b=4$. Compute $A=g^a$, $B=g^b$, and the shared secret $g^{ab}$, reducing the
exponent modulo the group order $22$. Verify Alice's and Bob's computations
agree.
\end{exercise}

\begin{exercise}
Explain precisely why a server that caches one ``ephemeral'' DH value $g^b$ and
reuses it for an hour of connections \emph{still authenticates correctly} but
\emph{loses forward secrecy}. What does an attacker who steals that cached $b$
recover, and over what set of sessions?
\end{exercise}

\begin{exercise}
Consider the signed handshake but with Bob signing only $B$, not $A\concat B$.
Give a concrete replay attack in which Eve causes Alice to accept a session using
a $B$ that Bob produced in an \emph{earlier, unrelated} session. Then show that
signing $A\concat B$ defeats your attack.
\end{exercise}

\begin{exercise}
Prove a small lemma about HKDF's domain separation under the assumption that
HMAC keyed by $\mathrm{PRK}$ is a pseudorandom function. Argue that for distinct
labels $\mathrm{info}_1\neq\mathrm{info}_2$, an adversary given the expanded
output for $\mathrm{info}_1$ cannot distinguish the output for $\mathrm{info}_2$
from random. (Hint: reduce to the PRF security of HMAC; the labels become
distinct inputs.)
\end{exercise}

\begin{exercise}
A protocol designer proposes to skip the Extract stage of HKDF and feed the raw
Diffie-Hellman element $g^{ab}\in\Zns{p}$ directly into Expand as the
$\mathrm{PRK}$. Give two distinct reasons this is dangerous, one about the
\emph{non-uniformity} of $g^{ab}$ and one about \emph{small-subgroup} structure
(what if $g^{ab}$ lands in a small subgroup?). Relate each to a property HKDF's
Extract is supposed to provide.
\end{exercise}

\begin{exercise}
Reconstruct the \textsc{Logjam}-style downgrade at the protocol level. Alice
offers suites $\{\text{ECDHE-strong},\ \text{DHE-export-512}\}$; Eve wants the
export suite chosen. Write the message Eve forwards to Bob, the message Bob
returns, and then show exactly where the final transcript-MAC check (if present)
would have caught her --- and why it does not if the handshake omits transcript
authentication.
\end{exercise}

\begin{exercise}
In TLS 1.3 0-RTT, explain why data in the client's first flight is replayable
while data sent \emph{after} the server's ServerHello is not. Identify the
specific ingredient (from this chapter) that the later data has and the 0-RTT
data lacks. Then classify which of these HTTP requests are safe to send as 0-RTT:
\texttt{GET /balance}, \texttt{POST /transfer}, \texttt{GET /logo.png},
\texttt{POST /search}.
\end{exercise}

\begin{exercise}
In Signal's X3DH secret, identify which of the four DH terms would still be
computable by an attacker who has stolen \emph{Bob's long-term identity private
key but none of his ephemeral/prekey secrets}, and which remain out of reach.
Use this to argue that X3DH gives forward secrecy for the initial message with
respect to identity-key compromise.
\end{exercise}

\begin{exercise}
Mutual authentication: modify the signed-ephemeral-DH protocol so that Bob also
verifies Alice. State exactly what Alice signs and what Bob checks, and explain
why both signatures must cover the \emph{same} transcript value rather than each
party signing only its own ephemeral.
\end{exercise}
